% ============================================================
% 第 2 章 数据预处理
% ============================================================

\section{数据预处理}
\label{sec:preprocess}

\subsection{符号约定}

为避免各问各自造符号，本文全文只使用一套符号，首次出现时定义，后文直接复用；角标最多两级（如 $x_{i,j}$），不引入三层及以上的上下标。主要符号见表~\ref{tab:notation}。

\begin{longtable}{@{}p{0.19\textwidth} p{0.49\textwidth} p{0.24\textwidth}@{}}
  \caption{全文符号约定}\\
  \toprule
  符号 & 含义 & 维度/取值 \\
  \midrule
  \endhead
  $N,P,S,G$ & 样本数、化合物数、感官维度数、风格类别数 & 正整数 \\
  $i,j,s,g$ & 样本、化合物、感官维度、风格索引 & $i=1,\dots,N$；其余同理 \\
  $k$ & 感官等级索引 & $k=0,1,\dots,6$ \\
  $\vx_i$ & 第 $i$ 个样本的化合物特征向量 & $\R^{P}$ \\
  $x_{i,j}$ & 第 $i$ 个样本第 $j$ 个化合物的特征值（逐列 log10 或线性，见变换表） & $\R$ \\
  $\bX$ & 样本—化合物特征矩阵 & $\R^{N\times P}$ \\
  $r_{j,j'}$ & 化合物 $j$ 与 $j'$ 的 Pearson 相关系数 & $[-1,1]$ \\
  $\bR$ & 化合物相关矩阵 & $\R^{P\times P}$ \\
  $a_{j,j'}$ & 共表达图邻接矩阵元素 & $\{0,1\}$ \\
  $\bA$ & 共表达图邻接矩阵 & $\{0,1\}^{P\times P}$ \\
  $y_{i,s}$ & 第 $i$ 个样本第 $s$ 个感官维度的评分 & $\{0,\dots,6\}$ 或连续潜变量 \\
  $\vy_i$ & 第 $i$ 个样本的感官评分向量 & $\R^{S}$ \\
  $\bY$ & 样本—感官评分矩阵 & $\R^{N\times S}$ \\
  $z_i$ & 第 $i$ 个样本的风格标签 & $\{1,\dots,G\}$ \\
  $\Dcal_g$ & 属于风格 $g$ 的样本集合 & $\Dcal_g\subseteq\{1,\dots,N\}$ \\
  $\bbeta_s$ & 第 $s$ 个感官维度的回归系数向量 & $\R^{P}$ \\
  $\bB$ & 各感官维度回归系数组成的矩阵 & $\R^{P\times S}$ \\
  $w_{j,s}$ & 化合物 $j$ 对感官维度 $s$ 的重要性 & $\R$ \\
  $\hat{y}_{i,s}$ & 模型预测的感官评分 & $\R$ \\
  $\tilde{\vx}$ & 生成的化合物谱 & $\R^{P}$ \\
  $\kappa_{s,k}$ & 感官维度 $s$ 的第 $k$ 个有序阈值 & $\R$ \\
  $\tau_j$ & 化合物 $j$ 的缺陷阈值 & $\R$ \\
  $q_{g,j}$ & 风格 $g$ 中化合物 $j$ 的 75\% 分位数 & $\R$ \\
  $\pi_{j,s}$ & 化合物 $j$ 在感官维度 $s$ 的 bootstrap 入选频率 & $[0,1]$ \\
  $\btheta$ & 模型参数 & — \\
  $\Lcal$ & 损失函数 & — \\
  \bottomrule
\end{longtable}
\label{tab:notation}

\subsection{数据概况与表间对齐}

附件各表的样本标识并不完全一致。\texttt{1.csv}、\texttt{4.csv}、\texttt{7.csv} 均含 \texttt{beer} 列，取值为 1--250 的唯一编号；\texttt{4.csv} 与 \texttt{7.csv} 的 \texttt{beer\_id} 存在大量科学计数法存储，无法作为可靠主键。因此本文以 \texttt{beer} 为唯一连接键，先建立样本主表，再逐表左连接，连接后核对每列的非空数量与取值范围。

样本主表共 250 行、231 个化合物特征。自然缺失 73 个单元（0.12\%），主要集中在 \texttt{amyl\_acetate}（27）、\texttt{aroma\_hops\_sum}（13）、\texttt{methylpyrazine}（12）、\texttt{2,3-butanedione}（9）与 \texttt{3-methylbutanal}（7）。感官目标共 50 维，列标准化后的标准差为 0.37--0.86，目标间相关中位数为 0.016，但存在若干强相关对（最高 $r=0.98$），说明不能把 50 个维度当作完全独立的任务。风格标签共 22 类，最大类 Blond 31 个样本、Tripel 28 个，最小类 Faro 3 个；小样本风格在问题三中采用收缩估计与 bootstrap 报告不确定性。

化合物丰度的变换并不是整表统一的。将 \texttt{8.xlsx} 的 Calibrated chemical dataset 与 \texttt{1.csv} 的 70 个公共化合物列逐列比较：37 列满足 $\text{calibrated}=10^{x_{i,j}}$（log10），33 列满足 $\text{calibrated}\approx x_{i,j}$（线性浓度）；其余 161 列没有校准浓度，按负值与值域判定为 log10。本文据此建立逐列变换表，只在解释浓度时按列还原，建模时统一用分位数变换，避免把 $10^{x}$ 误套到线性列上。完整的 231 列变换表随支撑材料提供，正文只报告 37/33/161 的统计与典型列（表~\ref{tab:transform-summary}）。

\begin{table}[htbp]
  \centering
  \caption{\texttt{1.csv} 逐列变换核对结果}
  \label{tab:transform-summary}
  \small
  \begin{tabular}{@{}l c l l@{}}
    \toprule
    类型 & 列数 & 判定依据 & 典型列 \\
    \midrule
    校准 log10 & 37 & $\text{calibrated}=10^{x}$ & \texttt{ethyl\_acetate}、\texttt{DMS}、\texttt{pH} \\
    校准线性 & 33 & $\text{calibrated}\approx x$ & \texttt{linalool}、\texttt{geraniol}、\texttt{2,3-butanedione} \\
    推断 log10 & 161 & 负值或值域 $<12$ & \texttt{X4VG}、萜类、酯类 \\
    \bottomrule
  \end{tabular}
\end{table}

感官数据方面，\texttt{3.xlsx} 只提供评分卡结构与 0--6 档位，不提供已填写的分数；\texttt{4.csv} 与 \texttt{7.csv} 给出的是标准化/中心化后的连续值，含负值与缺失。两者不能直接当作 0--6 标签使用。本文的处理是：以连续值作为监督信号训练回归或有序模型，再通过稳健区间映射与分位数校准得到 0--6 分级；所有涉及“真实等级”的指标仅在模拟等级或专家复核数据上报告，并在文中注明。

\texttt{4.csv}（训练面板）与 \texttt{7.csv}（Ratebeer 描述符）并非同一套感官体系。对可对应的概念逐对计算相关：酸感 0.824、酒精感 0.653、苦味 0.653、甜感 0.573，说明两者在主要口感维度上外部一致；但 DMS 仅 0.119、反式-2-壬烯醛 0.194、双乙酰 $-0.009$，说明痕量缺陷维度不能靠 \texttt{7.csv} 校验。本文因此把 \texttt{7.csv} 用作主要口感维度的外部稳健性检查，而不把它当作缺陷标签。

进一步用 Kaggle 的 22 个风格原型做独立校验：它与 \texttt{7.csv} 的风格均值在酸感（$r=0.890$）、麦芽（0.798）、酒花（0.819）、苦味（0.748）、酒精（0.688）、甜感（0.697）上方向一致；果味（0.147）与酒体（$-0.267$）则不可靠。与 \texttt{4.csv} 的相关整体更弱（酒精 0.262、麦芽 0.429、酒花 0.449），说明训练面板与评论词频并不是同一套感官体系。外部数据只用于风格级排序校验与先验，不能当作个体级 0--6 标签；这一限制在第~\ref{sec:q3} 节生成评价中同样适用。

\subsection{缺失值与异常值处理}

\texttt{1.csv} 与 \texttt{7.csv} 均存在缺失。对化合物矩阵，本文用缺失掩码 $\bM\in\{0,1\}^{N\times P}$ 记录观测位置，$m_{i,j}=1$ 表示第 $i$ 个样本第 $j$ 个化合物被观测到。缺失处理遵循两条原则：其一，任何插补、标准化与阈值选择都只在训练折内拟合，验证/测试折只做变换，防止信息泄漏；其二，插补方法作为敏感性分析的一个因素，而不是只保留一种“最优”结果。

基线插补采用训练集逐列中位数；Q1 另加缺失指示特征，Q3 对 \texttt{X\_cal} 的 3 个负值截断为 0（\texttt{betaglucan.mg.L.} 与 \texttt{unfermentables}）。异常值未做统一删除：化合物侧的极端值多为真实风味差异，删除会损失信息；本文在 Q1 的分位数变换、Q3 的 log1p 与分位数校准中降低其影响，并在结果中保留极端样本。缺失数据的统计处理框架参照 Little 与 Rubin 的讨论\citep{little2019statistical}。

\subsection{变换、标准化与特征构造}

化合物特征以 log10 为主，但 33 个校准列保留线性浓度；逐列变换后，不同化合物量纲差异仍然很大，本文对特征做训练集内标准化：
\begin{equation}
  \tilde{x}_{i,j}=\frac{x_{i,j}-\mu_j}{\sigma_j},
  \label{eq:standardize}
\end{equation}
其中 $\mu_j$、$\sigma_j$ 为训练集第 $j$ 列的均值与标准差。对组成型数据，本文另做中心对数比（CLR）变换作为敏感性分析，检查结论是否依赖单一变换。

特征构造遵循“可解释优先”的原则：保留原始化合物列作为主特征；按 \texttt{8.xlsx} 的 Sensory names 与化合物类别，把化合物归入酒花萜类、酵母酯类、高级醇类、醛酮类、麦芽来源与其他六组，构造组内均值与总和特征；对缺陷化合物单独建组，便于问题三的约束与问题二的合理性解释。所有派生特征在论文中注明来源列与计算方式，不使用无法追溯的临时特征。

\subsection{外部数据与补充信息}

题目给出的 Kaggle Top Beer Information 数据集\citep{kaggle2026topbeer} 已实际下载并核对：主表 5,558 行 $\times$ 21 列，来源为 BeerAdvocate 评论词频，许可为 CC BY 4.0。它与附件 250 款酒没有共同的酒名、酒厂或条码字段，无法逐酒连接；本文只在 22 个 \texttt{tasting\_category\_fine} 风格层级使用它：将 112 个 Kaggle 风格中的 96 个映射到附件 22 个风格，覆盖 4,760/5,558 行（85.6\%）。具体用途有三：其一，为 Q1 提供风格级感官先验与外部稳健性校验；其二，为 Q3 提供 22 个风格的感官原型与迁移方向；其三，用风格级 ABV 做交叉检查（$r=0.855$），而 IBU 的相关仅 0.449，故不把 IBU 强行用于个体样本。外部数据只做风格级辅助，不替代 \texttt{1.csv}/\texttt{2.csv} 的化合物建模。

\subsection{预处理流程与复现}

预处理流程如算法~\ref{alg:preprocess} 所示。该流程只依赖附件原始数据与本文代码，不依赖任何手工修改的中间表；随机种子、缺失率与折数统一记录在配置文件与附录中。

\begin{algorithm}[htbp]
  \small
  \caption{多源数据对齐与预处理}
  \label{alg:preprocess}
  \Input{附件 \texttt{1.csv}--\texttt{8.xlsx}}
  \Output{对齐特征矩阵 $\bX$、评分矩阵 $\bY$、风格标签 $\{z_i\}$、缺失掩码 $\bM$}
  读取各表，统一以 \texttt{beer} 为样本主键，核对行数与唯一性\;
  用 \texttt{8.xlsx} 的校准化学数据逐列核验 \texttt{1.csv} 的变换（37 列 log10、33 列线性）\;
  建立缺失掩码 $\bM$，标记 \texttt{1.csv} 与 \texttt{7.csv} 的缺失位置\;
  按样本分层划分训练/验证/测试集，固定随机种子\;
  \For{每个训练折}{
    在训练折内拟合中位数插补、均值与标准差，并变换其余各折\;
  }
  解析 \texttt{7.csv} 的 \texttt{tasting\_category\_fine}，建立风格集合 $\Dcal_g$\;
  输出对齐后的数据与一份可复查的预处理日志\;
\end{algorithm}

\figauto{fig_preprocess.pdf}{数据预处理与表间对齐流程}{fig:preprocess}

预处理流程如图~\ref{fig:preprocess} 所示：先对齐样本主键，再核验变换、建立掩码、划分数据与拟合变换参数，最后输出可复查的日志。该顺序保证任何监督信息都不会从验证/测试折泄漏到训练折。

需要强调的是，本文不把 \texttt{4.csv} 与 \texttt{7.csv} 的标准化值反标准化后直接当作 0--6 分数：标准化是在不同样本集上分别完成的，反变换会引入不可识别的一致性问题。正确做法是保留连续监督信号，另建从连续潜变量到 0--6 等级的有序校准层，详见第~\ref{sec:q1} 节。
